With the two proposed algorithms defined we now detail how we implemented them and the experiments we ran to test their effectiveness.


\subsection{Algorithm Implementations}\label{sec:impl}

The client-side algorithm is implemented on top of the OpenAI API~\cite{noauthor_openaiopenai-python_2025} using its async client.
Whenever a chat completion or response is posted to the main model, we also post a chat completion/response to the speculative API endpoint asynchronously.
If using multiple speculative samples, then we send multiple asynchronous requests.
When the speculative model API calls return, we launch the speculated tools asynchronously and put their futures into a tool cache keyed on function name and arguments.
This is implemented as a custom Python interface using Python's built-in asyncio library.
When the main model returns we check the tool cache if it called a tool.
If there is a cache hit, then we wait on the future and use its result as the tool output.
Otherwise we call the tool as normal.

For the engine-side algorithm we modify a custom fork of the popular vLLM framework~\cite{kwon2023efficient} to implement our algorithm.
A tool cache endpoint is implemented as a standard HTTP POST API endpoint (see later API spec). Similar to the client-side setup, one or more speculative model endpoints run asynchronously in parallel with the main model (\emph{O1}). The results of the speculated tool calls are posted to the tool-cache endpoint. 
Building on vLLM’s speculative decoding infrastructure, we implement a tool proposer responsible for detecting tool-call boundaries in the generated token stream. When the start of a tool call is detected, the proposer drafts the tool call tokens matching the latest entry in the cache with the same tool name. These draft tokens are then validated using standard speculative sampling~\cite{leviathan_fast_2023}, maintaining correctness (\emph{O2}). When the end of a tool call is detected, the proposer performs a lookup in the cache using both the tool name and its arguments. On a cache hit, the corresponding tool-output tokens are drafted; in this case, all drafted tokens can be accepted directly, since the match is exact on both name and arguments (\emph{O3}). If the lookup results in a cache miss, generation stops, and the engine returns the tool call to the client for normal execution.
For the engine-side algorithm to work in multi-turn scenarios we ensure streaming is enabled, so the client knows when to start speculating new tools.

The custom API endpoint we use is detailed below. This simple interface is all that is needed for existing commercial endpoint providers to expose this optimization to users.
Providers are incentivized to use the \texttt{cache-tool-output} API as it can cut down on inference time and increase throughput.
End-user can be incentivized by reduced costs and faster turnaround times. The cost reduction is not necessarily implicit, however,
many providers already offer discounts on prompts that hit the prefix-cache and our proposed optimization would fit nicely within that existing pricing model.

\begin{apispecbox}{\,\texttt{POST} \small \texttt{ /cache-tool-output/\{response\_id\}}}

\textbf{Description.}  
Cache tool outputs that the inference engine can utilize to keep sequences resident after tool calls. 
Each entry is keyed by \texttt{name} and (optionally) \texttt{params} and is matched against future tool calls by the model in this response.

\medskip
\noindent\textbf{Request body.}  
JSON array of objects with the following fields:
\medskip

\small
\begin{tabularx}{\linewidth}{@{}p{0.26\linewidth} L@{}}
\toprule
\textbf{Field} & \textbf{Description} \\
\midrule
\texttt{name} &
(string, required) Unique tool name used as part of the cache key. \\[0.25em]

\texttt{params} &
(object, optional) Canonicalized tool parameters included in the cache key.
Calls with the same \texttt{name} and equivalent \texttt{params} share outputs. \\[0.25em]

\texttt{output} &
(any, required) Serialized tool output to be cached and reused when the model
emits a matching tool call. \\[0.25em]

\texttt{keep\_alive} &
(number, optional) Optional time-to-live for the cache entry (e.g., in seconds). \\
\bottomrule
\end{tabularx}
\normalsize

\medskip
\noindent\textbf{Response.}
{\small
\begin{verbatim}
200 OK
{
  "cached": <int>  // number of entries accepted
}
\end{verbatim}
}

\end{apispecbox}



\subsection{Testing Environment}\label{sec:test-env}

We ran our implementations and experiments on a single compute node with four 80GB NVIDIA A100 GPUs controlled by one AMD EPYC 7763 CPU with 64 physical cores each at 2.45 GHz.
Each GPU pair is connected by third generation NVLink with 25 GB/s per direction links and each GPU is connected to the CPU via PCIe 4.0.
We run on a variation of SUSE Enterprise Linux Server 15 SP5 with Python 3.12 and CUDA 12.9.
In our experiments, we run two vLLM inference servers: one for the main model and one for the speculative model. 
The main model $\mathcal{M}$ uses a single A100 for serving, while the speculative server hosts three data parallel instances of the speculative model $\mathcal{S}$ across the other three A100s.

\subsection{Experiments}\label{sec:experiments}

To evaluate our inference optimizations, we use the dataset of prompts and tools from the BFCL project~\cite{patil_gorilla_2023}: a benchmark used to evaluate LMs on their tool calling capabilities.
BFCL defines a large number of tools that span a wide variety of tasks.
As a leaderboard, BFCL presents results comparing how well different LMs perform at calling the correct tools.
This provides a great starting point for selecting small, accurate speculation models.

From the BFCL leaderboards, we selected the xLAM models~\cite{zhang_xlam_2024} as ideal candidates for speculative models.
They are Llama 1B, 3B, and 8B fine-tunings on tool datasets to predict tool calls.
During testing we found the 8B model to predict around 80\% of tool calls from BFCL correctly, which is in line with the results they present in their paper.
Since these models are small and accurate at predicting tools, we utilize them as speculative tool calling models.



BFCL defines pairs of input prompts and sets of tools that can be invoked for each prompt, but it does not provide concrete tool implementations (only specifications). Our goal in this work is to study the inference- and systems-level behavior of agents with tools.%
We therefore treat tools as black boxes with specified interfaces and latencies. For each BFCL tool, we pre-compute a representative output using a mix of LM and human authored responses, and store these outputs in an in-memory cache. During experiments, whenever the model calls a tool, the cached tool output is returned and its runtime is controlled via manually configured tool latencies.

Crucially, our algorithms and measurements depend only on \emph{when} tools are called and \emph{how long} they take to run, not on the content of the tool outputs themselves. Precomputing outputs in this way preserves the model’s compute behavior while allowing us to systematically control tool latencies. As a consistency check, we manually inspected model generations across our experiments and found them to be reasonable and in line with typical outputs from the base model. We also ran a small experiment using simple, hand-implemented tools with real executions and observed no qualitative change in model behavior or quantitative change in the performance results, supporting the validity of our tool setup for system-level performance evaluation.


To turn BFCL into an agent workload, we treat each prompt and toolset pair as an independent task for an asynchronous agent. We run $M$ concurrent agents with $M \in \{1, 8, 32\}$. 
Each agent samples a prompt and its associated tools. It will loop continually calling tools (when requested) and generating text until it is finished and produces a final response. Then it will receive another prompt and tools, and keep repeating this process with new prompts.
We run each experiment long enough for every agent to complete 32 such tasks, and record the performance metrics reported in Section~\ref{sec:metrics} over the full run.
During our experiments we use the gpt-oss-120b~\cite{openai_gpt-oss-120b_2025} model as our primary agent as it is a state-of-the-art reasoning model that might be commonly used in an agent.

We systematically sweep the main experimental hyperparameters to study their impact.
For tool behavior, we sample tool latencies from random normal distributions of two types:
\emph{short-latency} distributions with mean latencies $\in 
\{0, 0.1,\ldots,0.5\}$,  and \emph{long-latency} distributions with means $\in \{0,0.5,\ldots,3\}$.
For speculation, we vary the number of speculative samples per turn ($1, 3, 5, 7,$ and $9$) and the speculation model (xLAM 1B, 3B, and 8B).
For each point in this search space, we run the workload under three inference configurations: standard inference (no speculation), client-side speculation, and engine-side speculation.
Each of these experiments is run five times to account for variation and system noise.
Additionally, we run a subset of experiments with the client-side algorithm and commercial models gpt-5 and gpt-5-nano to study how well this approach works for commercial models.

